% !TEX root = ../../cookbook.tex
\newpage
\recipe{Pasta with Anchovies}{Original recipe: \url{ricette.giallozafferano.it}. Serves 4}{}{}

\subsection*{Ingredients}
\begin{ingredients}
	\item 400 g pasta (ideally spaghetti)
	\item 30 g anchovies in oil, drained
	\item 35 g tomato paste
	\item 30 g pine nuts
	\item 1 clove of garlic
	\item 1 small bunch of parsley
	\item \sfrac{1}{2} fresh chili pepper
	\item White wine for deglazing
	\item Extra virgin olive oil, as needed
	\item 1 tablespoon of mascarpone per person (optional)
\end{ingredients}

% --- Preparation ---
Boil the \textbf{pasta}, setting aside a small glass of the cooking water.
Divide the \textbf{anchovies} into thirds, pouring their oil directly into a pan to heat. 
Melt \sfrac{1}{3} of the anchovies in the pan over medium heat. 
Meanwhile, on a cutting board, finely chop the rest of the anchovies together with the \textbf{pine nuts}, the \textbf{parsley}, the \textbf{garlic}, and the \textbf{chili pepper}.

Once the anchovies in the pan have melted, add a little cooking water and the \textbf{tomato paste}.
Let it reduce, then add the aromatic mixture and sauté it for a couple of minutes\footnote{If necessary, add a little more olive oil.}. 
Deglaze with the \textbf{white wine}. 
Add the al dente pasta, drained, to the pan, along with a bit more of the cooking water if needed. 
If desired, add a tablespoon of fresh mascarpone per person. 

% --- Description ---
\begin{center}
	\begin{minipage}{\textwidth} % Expanded to full text width for 3 columns
		\begin{recipeblock}
			\begin{multicols}{3}
				\footnotesize
				\centering % Centers text inside each column
				\setlength{\parskip}{1em} % Adds space between stanzas
				
				Acciughe \\
				Acciughe salate. \\
				Acciughe sebbele \\
				Trote salmonate. 
				
				Salmone \\
				Capparossoli,\\
				Risotto \\
				Al nero di seppia. 
				
				Spaghetti \\
				All'aragosta. 
				
				Anguilla \\
				Frittura mista, \\
				Polpa di granchio \\
				Scampi alla griglia. 
				
				Salmone \\
				Capparossoli, \\
				Risotto \\
				Al nero di seppia. 
				
				Spaghetti \\
				All'aragosta. 
				
				Spaghetti allo scoglio \\
				Pizza mari e monti. \\
				Paste di gamberi, 
				
				Ma soprattutto \\
				Acciughe \\
				Acciughe salate. 
				
				Peoci e trote, \\
				Salmone e seppie, \\
				Anguille, granchio, \\
				Scampi, salmone. 
				
				Acciughe \\
				Acciughe salate. \\
				\vspace{\baselineskip}
				\normalsize \textit{---Superlova}
			\end{multicols}
		\end{recipeblock}
	\end{minipage}
\end{center}